\begin{pastexam}[label={px:ncku-114-ee-ds-01}]{成大電機所 114 學年度 單選題第 1 題}
（單選；每題 5 分，答錯倒扣 2 分）

A is a two dimension array, the locations of A(3, 2) and A(2, 3) are 1115 and 1207,
respectively. Assume that each element occupies two addresses. Please give the location of
A(5, 10).
\begin{enumerate}[label=\alph*)]
\item 1841
\item 1851
\item 1861
\item 1871
\item None of the above.
\end{enumerate}
\end{pastexam}

\begin{pastsoln}
答案是 d) 1871。

題目沒說是以列為主（row-major）還是以行為主（column-major），也沒給列數、行數，要從兩個已知
位址推回來（資結書第~4.1~節）。 設陣列有 $R$ 列、$C$ 行，每個元素佔 $2$ 個位址。

\textbf{以列為主？} 同一列相鄰兩格差 $2$，下一列差 $2C$。 $A(2, 3)$ 在 $A(3, 2)$ 的前一列、
後一行，位址差 $\mathrm{loc}(3, 2) - \mathrm{loc}(2, 3) = 2C - 2 = 1115 - 1207 = -92$，得
$C = -45$，不可能。

\textbf{以行為主。} 同一行相鄰兩格差 $2$，下一行差 $2R$。
\[
\mathrm{loc}(2, 3) - \mathrm{loc}(3, 2) = 2R \cdot (3 - 2) + 2 \cdot (2 - 3) = 2R - 2 = 1207 - 1115 = 92,
\]
得 $R = 47$。 於是
\[
\mathrm{loc}(5, 10) = \mathrm{loc}(3, 2) + 2\bigl(R \cdot (10 - 2) + (5 - 3)\bigr)
= 1115 + 2(47 \cdot 8 + 2) = 1115 + 756 = 1871.
\]
用另一個已知點驗算：$\mathrm{loc}(2, 3) + 2(47 \cdot 7 + 3) = 1207 + 664 = 1871$，一樣。 這個
算法不需要知道索引的下界是 $0$ 還是 $1$，只用到兩點的差。

其他選項：a) 1841、b) 1851、c) 1861 分別少了 $30, 20, 10$ 個位址，都對不上 $R = 47$；有正確答案，
所以不是 e)。
\end{pastsoln}
